% Nuclear Stopping Cross Section and Straggling

\subsubsection{Nuclear stopping cross section}
%%%%%%%%%%%%%%%%%%%%%%%%%%%%%%%%%%%%%%%%%%%%%%
%
When a projectile travels a distance $\mathrm{d} r$ in a target with randomly arranged atoms of atomic density $N$ (units: atoms/\AA$^3$), it passes by $2 \pi p \, \mathrm{d} p \, \mathrm{d} r N$ atoms at distances between $p$ and $p + \mathrm{d} p$. According to \cref{eq:sn_physics_dE}, in a collision with impact parameter $p$, the projectile transfers the energy
%
\begin{equation}
    \label{eq:sn_sn_dE}
    T_\mathrm{n} = \gamma \, E \, \sin^2 \frac{\theta}{2} ~, \quad
    \gamma = \frac{ 4 M_1 M_2 }{ (M_1 + M_2)^2 }
\end{equation}
%
to the target atom, where $\theta = \theta(E, p)$ can be calculated as described in \cref{s:cmangle}. The average total energy loss $- \overline{\mathrm{d} E}$ of the projectile to target nuclei while traveling a distance $\mathrm{d} r$ therefore is
%
\begin{equation}
    - \overline{\mathrm{d} E} = \int_0^\infty T_\mathrm{n}(E,p) \,
        2 \pi p \, \mathrm{d} p \, \mathrm{d} r N
\end{equation}
%
The average energy loss to nuclei per path length (also called ``stopping force'') therefore is
%
\begin{equation}
    \label{eq:sn_sndef}
    - \overline{\left( \frac{ \mathrm{d} E }{ \mathrm{d} r } \right)}_\mathrm{n} =
        N \, S_\mathrm{n}(E)
\end{equation}
%
where the stopping cross section (also called ``stopping power'') is given by
%
\begin{equation}
    S_\mathrm{n}(E) = \pi \int_0^\infty T_\mathrm{n}(E,p) \, 2p \, \mathrm{d} p =
        \pi \gamma E \int_0^\infty \sin^2 \frac{\theta}{2} \, 2p \, \mathrm{d} p ~.
\end{equation}
%
In terms of the reduced quantities Eqs.~(\ref{eq:sn_reduced_lengths}) and (\ref{eq:sn_reduced_energy}) the stopping cross section may be written
%
\begin{equation}
    \label{eq:sn_sn1}
    S_\mathrm{n}(E) = \pi \gamma (E/\varepsilon) a_\mathrm{I}^2 \, s_\mathrm{n}(\varepsilon)
\end{equation}
%
with the reduced nuclear stopping cross section
%
\begin{equation}
    s_\mathrm{n}(\varepsilon) = \varepsilon \int_0^\infty \sin^2 \frac{\theta(\varepsilon,P)}{2} \, 2 P \, \mathrm{d} P
\end{equation}
%
With Eqs.~(\ref{eq:sn_reduced_energy}) and (\ref{eq:sn_coulomb}), \cref{eq:sn_sn1} can be rewritten
%
\begin{equation}
    \label{eq:sn_sn}
    S_\mathrm{n}(E) =  \pi \frac{ 4 M_1 }{ M_1 + M_2 } Z_1 Z_2 (14.4 \,\mathrm{eV\mathAA}) \, a_\mathrm{I} \, s_\mathrm{n}(\varepsilon)
\end{equation}

\begin{figure}
    \centering
    \includegraphics[scale=0.5]{sn_p0.23_p1_homo.pdf} ~~
    \includegraphics[scale=0.5]{sn_p0.23_p1_homo_e0.01.pdf} ~~ \\[\baselineskip]
    \includegraphics[scale=0.5]{sn_p0.5_p0.67_homo.pdf} ~~
    \includegraphics[scale=0.5]{sn_p0.5_p0.67_homo_e0.01.pdf}
    \caption{Reduced nuclear stopping cross section for the NLHlin potential (colored lines) of homonuclear atom pairs ($Z_1=Z_2$) as a function of reduced energy $\varepsilon$ (left) and for $\varepsilon = 0.01$ as a function of atomic number (right).
        Two different screening lengths have been used.
        Upper row: ZBL screening length \cref{eq:sn_aZBL}; lower row: Firsov screening length \cref{eq:sn_aF}
        For comparison, the stopping cross sections of the ZBL potential (upper row) and the KrC potential (lower row) are depicted by black dashed lines.
        Note that these two potentials yield universal reduced stopping cross sections with the respective screening lengths.
        In the left panels the atomic numbers are color-coded.
        The dotted parts of the lines correspond to energies $< 10$~eV.
    }
    \label{fig:sn_sn}
\end{figure}
%
The reduced nuclear stopping cross sections $S_\mathrm{n}$ for the NLHlin potentials are shown as a function of reduced energy $\varepsilon$ in \cref{fig:sn_sn} for homo-nuclear pairs.
Considerable variation is observed, in particular at low energies, which cannot be described by universal interatomic potentials like ZBL and KrC.
Comparing the results for the universal and the NLHlin potentials at $\varepsilon < 1$, the ZBL potential lies at the upper limit of the NLHlin potentials, while the KrC potential better approximates the average of the heavier elements ($Z \gtrsim 40$).
On the other hand, KrC does even worse than ZBL for the lighter elements ($Z \lesssim 20$), which is mainly due to the Firsov screening length (check!).

\begin{figure}
    \centering
    \includegraphics[scale=0.5]{sn_ratio_homo_zbl_at_e.pdf} ~~
    \includegraphics[scale=0.5]{sn_ratio_homo_krc_at_e.pdf}
    \caption{Ratio of the stopping cross section calculated using the NLHlin potential and the ZBL potential (left) or KrC potential (right) as a function of atomic number for four different energies. The data has been obtain for the homonuclear case ($Z_1 = Z_2$).}
    \label{fig:sn_sn_ratio}
\end{figure}
%
For the light ions the low values of the reduced energy are of no relevance since they correspond to energies where the BCA breaks down.
It is therefore more instructive to plot the nuclear stopping cross section as a function of the actual energy.
\cref{fig:sn_sn_ratio} shows the ratios of the stopping cross sections calculated using the NLHlin potential and the ZBL potential or KrC potential as a function of atomic number for four different energies for the homonuclear case ($Z_1 = Z_2$).
Compared to when the stopping cross section is plotted for constant reduced energies (\cref{fig:sn_sn}), the deviations from 1 appear more evenly distributed between light and heavy ions.

It is noted that in BCA simulations, impact parameters must be limited by some maximum value $p_\mathrm{max}$.
This value is often chosen relatively small, so that considerable energy loss is neglected at low energies (check!).
Thus, the effective stopping cross section is lowered in the simulation, possibly towards the values for the NLHlin potential, however, in an unsystematic way.

\begin{figure}
    \centering
    \includegraphics[scale=0.5]{sn_zbl_all_e0.01.pdf} ~~
    \includegraphics[scale=0.5]{sn_krc_all_e0.01.pdf}
    \caption{Reduced nuclear stopping cross section for the NLHlin potential of all atom pairs as a function of atomic number $Z_1$ at a reduced energy of $\varepsilon = 0.01$.
    Two different screening lengths have been used.
    Left: ZBL screening length \cref{eq:sn_aZBL}; right: Firsov screening length \cref{eq:sn_aF}
    For comparison, the stopping cross sections of the ZBL potential (left) and the KrC potential (right) are depicted by black dashed lines.
    Note that these two potentials yield universal reduced stopping cross sections with the respective screening lengths.
    The atomic numbers $Z_2$ are color-coded.
    }
    \label{fig:sn_sn_all}
\end{figure}
%
The stopping cross sections for all atom pairs at $\varepsilon = 0.02$ are depicted in \cref{fig:sn_sn_all}.
The results are in line with those for the homo-nuclear pairs (\cref{fig:sn_sn}).


\subsubsection{Nuclear straggling}
%%%%%%%%%%%%%%%%%%%%%%%%%%%%%%%%%%
%
It may be shown \cite{Bohr1948} that the variance of the energy loss due to nuclear collisions over a path length $\mathrm{d} r$ in a target with randomly positioned atoms is given by
%
\begin{equation}
    \label{eq:sn_sn_Estraggle}
    \overline{ \left( \mathrm{d} E - \overline{\mathrm{d} E} \right)^2 } =
        \int_0^\infty T_\mathrm{n}^2(E,p) \,
        2 \pi p \, \mathrm{d} p \, \mathrm{d} r N =
    N \mathrm{d} r \, Q_\mathrm{n}(E)
\end{equation}
%
with
%
\begin{equation}
    Q_\mathrm{n}(E) = \int_0^\infty T_\mathrm{n}^2(E,p) \, 2 \pi p \, \mathrm{d} p ~,
\end{equation}
%
which is usually called ``nuclear straggling''. Using \cref{eq:sn_sn_dE}, we obtain
%
\begin{equation}
    Q_\mathrm{n}(E) = \pi \gamma^2 E^2 \int_0^\infty \sin ^4 \frac{\theta}{2} \, 2 p \, \mathrm{d} p ~.
\end{equation}
%
In terms of the reduced quantities Eqs.~(\ref{eq:sn_reduced_lengths}) and (\ref{eq:sn_reduced_energy}) the nuclear straggling may also be written
%
\begin{equation}
    \label{eq:sn_qn1}
    Q_\mathrm{n}(E) = \pi \gamma^2 (E/\varepsilon)^2 a_\mathrm{I}^2 \, q_\mathrm{n}(\varepsilon)
\end{equation}
%
with
%
\begin{equation}
    q_\mathrm{n}(\varepsilon) = \varepsilon^2 \int_0^\infty \sin^4 \frac{\theta(\varepsilon,P)}{2} \, 2 P \, \mathrm{d} P ~.
\end{equation}
%
With Eqs.~(\ref{eq:sn_reduced_energy}) and (\ref{eq:sn_coulomb}), \cref{eq:sn_qn1} can be rewritten
%
\begin{equation}
    \label{eq:sn_qn}
    Q_\mathrm{n}(E) =  \pi \left[ \frac{ 4 M_1 }{ M_1 + M_2 }
        Z_1 Z_2 (14.4 \,\mathrm{eV\mathAA}) \right]^2 \, q_\mathrm{n}(\varepsilon) ~.
\end{equation}

The results for the nuclear straggling for the NLHlin potentials are shown in \cref{fig:sn_qn}.
They compare similarly to the ZBL and KrC potential as the nuclear stopping cross sections.
%
\begin{figure}
    \centering
    \includegraphics[scale=0.5]{qn_p0.23_p1_homo.pdf} ~~
    \includegraphics[scale=0.5]{qn_p0.23_p1_homo_e0.01.pdf} \\[\baselineskip]
    \includegraphics[scale=0.5]{qn_p0.5_p0.67_homo.pdf} ~~
    \includegraphics[scale=0.5]{qn_p0.5_p0.67_homo_e0.01.pdf}
    \caption{Reduced nuclear straggling for the NLHlin potential (colored lines) of homonuclear atom pairs ($Z_1=Z_2$) as a function of reduced energy $\varepsilon$ (left) and for $\varepsilon = 0.01$ as a function of atomic number (right).
    Two different screening lengths have been used.
    Upper row: ZBL screening length \cref{eq:sn_aZBL}; lower row: Firsov screening length \cref{eq:sn_aF}
    For comparison, the stopping cross sections of the ZBL potential (upper row) and the KrC potential (lower row) are depicted by black dashed lines.
    Note that these two potentials yield universal reduced stragglings with the respective screening lengths.
    In the left panels the atomic numbers are color-coded.
    The dotted parts of the lines correspond to energies $< 10$~eV.
    }
    \label{fig:sn_qn}
\end{figure}

\subsubsection{Analytic expressions}
%%%%%%%%%%%%%%%%%%%%%%%%%%%%%%%%%%%%
%
The results for the nuclear stopping cross section and straggling can be fitted by analytical expressions within a few percent as proposed in \cite{ZBL85}. For the stopping cross section we use
%
\begin{equation}
    \label{eq:sn_sn_fit}
    s_\mathrm{n}(\varepsilon) = \frac{ \ln(1 + a \, \varepsilon) }{
                          2 ( \varepsilon + b \, \varepsilon^c + d \, \varepsilon^{0.5})}
\end{equation}
%
In \cite{ZBL85} it is suggested to use $s_\mathrm{n} = \ln \varepsilon / 2 \varepsilon$ for $\varepsilon > 30$, which is referred to as ``unscreened nuclear stopping'' without further explanation.
We did not find it necessary to apply this special treatment for high energies, and found actually better fitting of the data using the expression \cref{eq:sn_sn_fit} over the whole energy range (at least) up to $\varepsilon = 10^4$.

For the nuclear straggling we use \cite{ZBL85}
%
\begin{equation}
    \label{eq:sn_qn_fit}
    q_\mathrm{n}(\varepsilon) = \frac{ 1 }{ 4 + a \, \varepsilon^b + c \, \varepsilon^d }
\end{equation}
%
The values $a$, $b$, $c$, and $d$ have been fitted to energies between $\varepsilon = 3 \times 10^{-5}$ and $10^4$ for both nuclear stopping cross section and nuclear straggling for all atom pairs.

\begin{figure}
    \centering
    \includegraphics[scale=0.5]{sn_fit_nlhlin_B_Si.pdf} ~~
    \includegraphics[scale=0.5]{qn_fit_nlhlin_B_Si.pdf}
    \caption{Data and fits for $Z_1 = 5$, $Z_2 = 14$.}
    \label{fig:sn_fit_Z5_14}
\end{figure}
%
\begin{figure}
    \centering
    \includegraphics[scale=0.5]{sn_fit_error_nlhlin.pdf} ~~
    \includegraphics[scale=0.5]{qn_fit_error_nlhlin.pdf}
    \caption{Root-mean-square fit errors of the nuclear stopping cross section (left) and nuclear straggling (right) for all atom pairs.}
    \label{fig:sn_fit_error}
\end{figure}
%
\Cref{fig:sn_fit_Z5_14} shows, as an example the data and fits for $Z_1 = 5$, $Z_2 = 14$.
The root-mean-square fit errors for all atom pairs are shown in \cref{fig:sn_fit_error}.
Most of the fit errors is between 1 and 2.5\% for the nuclear stopping cross section and between 1.5 and 4\% for the nuclear straggling, with a mean of 1.76\% and 2.95\%,respectively, and a few errors up to 6\% and 8\%, respectively.